%%
%% HFF for Heterogeneous Objectives and Tournaments
%% Working draft — building from the appendix forward.
%%
\documentclass[sigconf]{acmart}

\usepackage{amsmath}
\usepackage{booktabs}
\usepackage{float}
\usepackage{algorithm}
\usepackage{algorithmic}
\usepackage{geometry}

\begin{document}

\title{Hyperspherical Fitness Functions for Heterogeneous Objectives and Tournaments}

\author{Andrew James Morgan}
\email{andrew@gamakon.ai}
\affiliation{%
  \institution{Gamakon Ltd.}
  \city{London}
  \country{United Kingdom}
}
\renewcommand{\shortauthors}{Morgan}

\maketitle

%% ============================================================================
%% BODY SECTIONS — stubbed, written last.
%% ============================================================================
\section{Introduction}
\label{sec:intro}
% TODO

\section{Method}
\label{sec:method}

\paragraph{Equivalence classes as tournaments.} An equality-saturation engine
(here egglog) takes an expression and a rule set and grows an e-graph until no
new equalities appear, producing an equivalence class: a set of forms that all
denote the same function. Equivalence holds in the \emph{goal} domain, so no
member is preferred on behaviour. Selecting one member --- \emph{extraction} ---
is therefore a tournament over goal-equivalent candidates, decided entirely on
properties of \emph{form}. Conventional extraction minimises a single scalar
cost; we instead score each member on a vector of form measures and rank by
hyperspherical fitness.

\paragraph{Pattern-dependent measures.} We introduce a library of guarded
measures written in an \texttt{awk}-like form
$\langle\textit{pattern}\rangle\ \{\textit{measure}\}$: a measure is collected
for a candidate only if its AST contains the pattern (Appendix~%
\ref{app:metrics}). Each measure is a cheap, deterministic, data-free tree walk,
defined to emit a value in $[0,1]$ with $0$ best. Because patterns fire
selectively, two members of the same class may report \emph{different} numbers
of measures, of different kinds --- a candidate with no trigonometric nodes
never reports the trig measures. The objective set is thus emergent from each
form's structure, not fixed across the class. We always include at least one
universally-firing measure --- node count (parsimony) --- so that every
candidate has a non-empty, non-zero measure vector and the projection below is
well-defined.

\paragraph{Per-class angles via the augmented projection.} For a candidate with
measure vector $\mathbf{x}\in[0,1]^k$ (its own $k$, set by which patterns fired),
we reuse the HFF-TrueNorth augmentation~\cite{morgan2026hff,stolfi1991}: the
direction $\mathbf{y}=\mathbf{x}/\|\mathbf{x}\|$ is undefined at the ideal
$\mathbf{x}=\mathbf{0}$, so we adjoin an energy score
$e=\max(0,\,1-\|\mathbf{x}\|^2/k)$ and form
$\mathbf{y}'=(\sqrt{1-e^2}\,\mathbf{y},\,e)\in S^{k}$, giving a single angle
\begin{equation}
\label{eq:class-angle}
\theta = \arccos(e), \qquad \theta=0 \ \text{iff all measures are } 0.
\end{equation}
Each equivalent form collapses to one angle: small magnitude (few, low penalties)
and balanced direction yield small $\theta$.

\paragraph{CDF correction makes the angles comparable.} A class of, say, 14 forms
yields 14 angles drawn from spheres of \emph{different} dimension $k$ --- the
counts are not directly comparable, because angular concentration depends on
$k$. We apply the CDF correction~\cite{cai2013angles}: under uniform sampling on
$S^{k-1}$,
\begin{equation}
\label{eq:class-cdf}
P(\theta\le t) = I_{\sin^2 t}\!\left(\tfrac{k-1}{2},\,\tfrac12\right),
\end{equation}
with $I_x(a,b)$ the regularised incomplete beta function. Each raw angle becomes
a dimension-adjusted percentile, so forms scored on 4 axes and on 7 axes rank on
one yardstick. Sorting the corrected percentiles totally orders the class, and
the minimum is the ``best'' equivalent form under the chosen patterns --- a
behaviour-preserving rewrite that is structurally preferred across all applicable
measures at once, with no hand-tuned cost weights.

\section{Experiments}
\label{sec:experiments}
% TODO: (1) GA tournaments (GECCO result). (2) e-graph extraction as a
% tournament over goal-equivalent forms, scored by the pattern-dependent
% metric library of Appendix~\ref{app:metrics}.

\section{Discussion}
\label{sec:discussion}
% TODO

\section{Conclusion}
\label{sec:conclusion}
% TODO

%% ============================================================================
%% APPENDIX — written first.
%% ============================================================================
\appendix

\section{Pattern-Dependent Metric Library}
\label{app:metrics}

The extraction tournament operates over a saturated equivalence class: every
member denotes the same function (equivalence holds in the \emph{goal} domain),
so no member is preferred on behaviour. Selection instead ranks members on
properties of \emph{form}. Each property is a cheap, deterministic, data-free
walk of the candidate's abstract syntax tree, and each is oriented so that
\textbf{0 is best} (lower is preferred), matching the minimisation target of
HFF-TrueNorth.

A metric is written in a guarded, \texttt{awk}-like form
\(\langle\textit{pattern}\rangle\ \{\textit{measure}\}\): the measure is
collected only for candidates whose AST contains the pattern. A candidate that
contains no trigonometric nodes simply does not report the trig metrics. This is
the source of \emph{heterogeneity}: two members of the same class may be scored
on different numbers of axes, of different kinds and units. The metrics are made
cross-comparable by the CDF-corrected angular distance of
Section~\ref{sec:method}, which converts a raw angle in a candidate's own
objective dimensionality into a dimension-adjusted percentile rank.

The library is open: any guarded measure that returns a non-negative scalar may
be added, and HFF assembles whichever subset applies per candidate. Table~%
\ref{tab:metrics} lists the initial set used in the extraction experiments.

Every measure is \textbf{defined to emit a value in $[0,1]$} (0 best), so a
candidate's measures form a clean vector with no per-batch rescaling: the vector
is L2-normalised, the angle to the pole is taken, and only the resulting angles
are CDF-corrected (each in its own dimensionality) to become cross-comparable.
The \emph{Normalisation} column gives the bounded form for each measure. Raw
unbounded counts use the saturating map $c \mapsto 1 - 1/(1+c/s)$ (0 at $c=0$,
$\to 1$ as $c$ grows, scale $s$ per measure); ratios and indicators are already
in $[0,1]$ by construction.

\begin{table*}[t]
\small
\centering
\caption{Initial pattern-dependent metric library. Each row is a guarded
measure $\langle\textit{pattern}\rangle\ \{\textit{measure}\}$ computed by one
AST walk; 0 is best, and the \emph{Normalisation} column maps it into $[0,1]$ at
source so the per-candidate measure vector needs no later rescaling. Self-nesting
of transcendentals is the strongest ``curve-fitter, not a law'' signal observed
in symbolic-regression near-misses.}
\label{tab:metrics}
\begin{tabular}{@{}llp{3.6cm}l@{}}
\toprule
\textbf{Pattern} & \textbf{Measure} & \textbf{Rationale} & \textbf{Normalisation $\to[0,1]$} \\
\midrule
\texttt{trig(trig(\ldots))} & \texttt{self\_nested\_trig\_count} & $\sin(\sin x)$: no physical law self-nests trig & $1-1/(1+c)$ \\
\texttt{f(f(\ldots))}, $f$ transc. & \texttt{self\_nested\_transc\_count} & $\exp(\exp)$, $\log(\log)$ — fitter signature & $1-1/(1+c)$ \\
\texttt{transc(transc(\ldots))} & \texttt{transc\_nesting\_depth} & any nonlinear-in-nonlinear depth & $1-1/(1+d)$ \\
\texttt{.*} & \texttt{max\_tree\_depth} & longest root-to-leaf path & $1-1/(1+d/s)$ \\
\texttt{.*} & \texttt{node\_count} & total nodes (parsimony) & $1-1/(1+n/s)$ \\
\texttt{Div(\_,0)|Log($\le$0)|Sqrt(neg)} & \texttt{domain\_violation\_count} & poles / domain breaches & $\mathbf{1}[c>0]$ \\
\texttt{Div(\_, var-expr)} & \texttt{variable\_denominator\_count} & asymptote / blow-up risk & $1-1/(1+c)$ \\
\texttt{Num} & \texttt{distinct\_constant\_count} & free parameters (overfit proxy) & $1-1/(1+c/s)$ \\
\texttt{transc op} & \texttt{transcendental\_count} & universal-approximator load & $1-1/(1+c/s)$ \\
\texttt{Pow(\_, non-int)} & \texttt{fractional\_power\_count} & radical noise & $1-1/(1+c)$ \\
\texttt{Neg|Sub} & \texttt{sign\_op\_count} & cancellation-prone structure & $1-1/(1+c/s)$ \\
\texttt{identical subtree} & \texttt{repeated\_subexpr\_count} & CSE-able $\Rightarrow$ bloated & $1-1/(1+c)$ \\
\texttt{.*} & \texttt{operator\_type\_diversity} & number of distinct op kinds & $k/k_{\max}$ \\
\texttt{Mul(Num,Mul(Num,\_))} & \texttt{unfolded\_constant\_count} & should have constant-folded & $1-1/(1+c)$ \\
\texttt{top-level Add chain} & \texttt{additive\_term\_count} & number of summands & $1-1/(1+c/s)$ \\
\texttt{.*} & \texttt{max\_arity\_used} & structural width & $a/a_{\max}$ \\
\texttt{Var} & \texttt{variable\_reuse\_count} & repeated-variable load & $1-1/(1+c/s)$ \\
\texttt{Pow(\_, large)} & \texttt{max\_exponent\_magnitude} & high-degree stiffness & $1-1/(1+|e|/s)$ \\
\texttt{.*} & \texttt{tree\_imbalance} & leaf-depth variance (ad-hoc shape) & $\tanh(\mathrm{var})$ \\
\texttt{.*} & \texttt{constant\_to\_variable\_ratio} & constant-heavy $\Rightarrow$ fitter & $n_c/(n_c{+}n_v)$ \\
\bottomrule
\end{tabular}
\end{table*}

\end{document}
